%!TEX TS-program = xelatex
%!TEX encoding = UTF-8 Unicode
%% bookspine.tex — NFU 工管系書背（獨立編譯）

\documentclass[a4paper,12pt]{article}
\usepackage{fontspec}
\usepackage{xeCJK}
\usepackage[a4paper,margin=1cm]{geometry}

%% 中文字型
\IfFontExistsTF{TW-Kai}{
  \setCJKmainfont{TW-Kai}[AutoFakeBold=2]
}{
  \IfFontExistsTF{AR PL KaitiM Big5}{
    \setCJKmainfont{AR PL KaitiM Big5}[AutoFakeBold=2]
  }{
    \setCJKmainfont{Noto Serif CJK TC}[AutoFakeBold=2]
  }
}

%% ------------------ 書背資料設定 ------------------------------

%% 系所：完整名稱
% 工業管理系工業工程與管理碩士班
\newcommand{\bsDeptVertical}{%
  \vbox to \bsTopNamesH{%
    \hbox to 1.2em{\hfil 工\hfil}\vfil
    \hbox to 1.2em{\hfil 業\hfil}\vfil
    \hbox to 1.2em{\hfil 管\hfil}\vfil
    \hbox to 1.2em{\hfil 理\hfil}\vfil
    \hbox to 1.2em{\hfil 系\hfil}\vfil
    \hbox to 1.2em{\hfil 工\hfil}\vfil
    \hbox to 1.2em{\hfil 業\hfil}\vfil
    \hbox to 1.2em{\hfil 工\hfil}\vfil
    \hbox to 1.2em{\hfil 程\hfil}\vfil
    \hbox to 1.2em{\hfil 與\hfil}\vfil
    \hbox to 1.2em{\hfil 管\hfil}\vfil
    \hbox to 1.2em{\hfil 理\hfil}\vfil
    \hbox to 1.2em{\hfil 碩\hfil}\vfil
    \hbox to 1.2em{\hfil 士\hfil}\vfil
    \hbox to 1.2em{\hfil 班\hfil}
  }%
}

% 國立虎尾科技大學
\newcommand{\bsUnivVertical}{%
  \vbox to \bsTopNamesH{%
    \hbox to 1.2em{\hfil 國\hfil}\vfil
    \hbox to 1.2em{\hfil 立\hfil}\vfil
    \hbox to 1.2em{\hfil 虎\hfil}\vfil
    \hbox to 1.2em{\hfil 尾\hfil}\vfil
    \hbox to 1.2em{\hfil 科\hfil}\vfil
    \hbox to 1.2em{\hfil 技\hfil}\vfil
    \hbox to 1.2em{\hfil 大\hfil}\vfil
    \hbox to 1.2em{\hfil 學\hfil}
  }%
}

%% 使用者修改資料
\def\bsYear{116}

\def\bsDegreeV{
  碩\\
  士\\
  論\\
  文
}

\def\bsTitleV{
  中\\
  文\\
  題\\
  目
}

\def\bsAuthorV{
  你\\
  的\\
  名\\
  字
}

%% ------------------ 書背尺寸設定 ------------------------------
%% A4 上呈現書背示意框
\newlength{\bsSpineH}
\setlength{\bsSpineH}{270mm}

\newlength{\bsSpineW}
\setlength{\bsSpineW}{20mm}

%% 系上附件十三：
%% 上方校名＋系所 6cm
\newlength{\bsTopNamesH}
\setlength{\bsTopNamesH}{6cm}

\newlength{\bsDegreeH}
\setlength{\bsDegreeH}{3cm}

%% 碩士論文 → 論文題目：1 cm
\newlength{\bsGapDegreeTitle}
\setlength{\bsGapDegreeTitle}{1cm}

%% 姓名區：2 cm
\newlength{\bsAuthorH}
\setlength{\bsAuthorH}{2cm}

%% 姓名 → 學年度：1 cm
\newlength{\bsGapAuthorYear}
\setlength{\bsGapAuthorYear}{1cm}

%% 學年度區：2 cm
\newlength{\bsYearH}
\setlength{\bsYearH}{2cm}

%% 書背底端：2 cm
\newlength{\bsBottomGap}
\setlength{\bsBottomGap}{2cm}

%% 垂直文字
\newcommand{\vtext}[1]{%
  \begin{tabular}[c]{@{}c@{}}
    #1
  \end{tabular}
}

\begin{document}
\thispagestyle{empty}
\begin{center}

  \setlength{\fboxsep}{0pt}
  \setlength{\fboxrule}{0.6pt}

  \framebox{%
    \begin{minipage}[t][\bsSpineH][t]{\bsSpineW}

      \centering
      \fontsize{12}{12}\selectfont
      \vspace{1.5cm}

      %% 1. 校名＋系所
    \begin{tabular}{@{}c@{\hspace{0.25cm}}c@{}}
        \bsDeptVertical &
        \bsUnivVertical
    \end{tabular}


      \begin{minipage}[c][\bsDegreeH][c]{\linewidth}
        \centering
        \vtext{\bsDegreeV}
      \end{minipage}


      %% 3. 空 1 cm
      \vspace*{\bsGapDegreeTitle}

      %% 4. 論文題目
      \vtext{\bsTitleV}\par
      \vfill

      %% 5. 姓名：2 cm 區域
      \begin{minipage}[c][\bsAuthorH][c]{\linewidth}
        \centering
        \vtext{\bsAuthorV}
      \end{minipage}

      %% 6. 空 1 cm
      \vspace*{\bsGapAuthorYear}

      %% 7. 學年度：2 cm 區域
      \begin{minipage}[c][\bsYearH][c]{\linewidth}
        \centering
        \vtext{
          \bsYear\\
          學\\
          年\\
          度
        }
      \end{minipage}

      %% 8. 底部留 2 cm
      \vspace*{\bsBottomGap}
    \end{minipage}%
  }
\end{center}
\end{document}
